\section{带权矩阵树定理与方向约定}\label{knowledge-matrix-tree-weighted}
实现见第~\ref{compact-MatrixTree}~节（第~\pageref{compact-MatrixTree}~页）。
返回的是每棵树的边权乘积，再对所有树求和，并对素数取模。
无权计数只需把所有边权设为 1。
重边是不同选择，自环不属于任何生成树；单点图的空树乘积为 1。
零权和负权都按域上的运算处理，因此结果为零可能来自抵消，不能作为连通性判定。

无向边 $uv$ 权值 $w$ 对拉普拉斯矩阵贡献为
$L_{uu},L_{vv}$ 各加 $w$，$L_{uv},L_{vu}$ 各减 $w$。
删除 root 行和列后求行列式，与 root 的选择无关。
对于每个点都能沿有向边走到根的树形图，
边 $u\to v$ 贡献 $L_{uu}\mathrel{+}=w$、$L_{uv}\mathrel{-}=w$；
删除根行列后得到内向树形图的权值和。
外向树形图先反转所有边，再使用同一构造。

例如只有一条边 $0\to1$ 时，内向根 1 的答案为该边权，内向根 0 为零；
外向根 0 的答案为该边权。
\href{https://www.luogu.com.cn/problem/P6178}{P6178} 的有向模式（type=1）指定以 1 为根的外向树；type=0 为无向图，
完整用法见第~\ref{usage-example-76}~节（第~\pageref{usage-example-76}~页）；驱动把点号减一，有向模式选择 \texttt{away\_from\_root}。
只直接构造 $(n-1)\times(n-1)$ 余子式，时间 $O(n^3+m)$、额外空间 $O(n^2)$。
行列式复用本书的素数模高斯消元，不支持将 mod 随意改成合数。

\subsection{为什么这个余子式恰好计数}
先证内向情形。删去根 $r$ 的坐标，令 $e_r=0$，其他 $e_i$ 为单位行向量。
余子式第 $i$ 行是 $\sum_{i\to j}w_{ij}(e_i-e_j)$。
按行的多线性展开，每项从每个非根点恰选一条出边，权为所选边权之积。
若这些边含有有向环，环上向量 $e_i-e_j$ 相加为零，该项行列式为零。
若无环，有限图中每条路径只能在没有出边的根终止。
按到根距离从小到大，同时重排行列，得到对角全为 1 的三角矩阵，行列式为 1。
因此只有内向生成树保留下来，每棵树的权恰计一次。

自环的向量为零；重边在展开时是不同选择；单点对应空行列式与空乘积 1。
无向边替换成权相同的两条反向弧，每棵无向生成树恰有一种朝根定向；
有向边全部反转则把外向树变成内向树。
整个论证是整数系数多项式恒等式，没有除法，故允许负权、零权及任意正模数。
素数限制来自消元实现，不来自矩阵树定理本身。

证明依据行列式多线性与有向树的出度刻画，参见
\href{https://ocw.mit.edu/courses/18-212-algebraic-combinatorics-spring-2019/resources/mit18_212s19_lec29/}{MIT 18.212 第29讲}；
这里直接按选边后的环与根路径给出主余子式的证明，不扩展为任意非主余子式公式。

\subsection{既有接口与验证记录}
测试独立枚举 $n-1$ 条边的所有选择，检查连通性或根向约束，
用任意精度整数计算边权乘积之和，再分别对多个素数取模。
两套实现覆盖所有根、三种方向、自环重边、负权及 64 位权值边界，
并通过 220 点完全图的加权 Cayley 公式检查。
P6178 完整驱动另测 300 点、十万条边以及输入输出约定。
普通与 ASan/UBSan 均通过。
两套代码已在 P6178 通过全部 11 个测试点：
\href{https://www.luogu.com.cn/record/297619886}{动态版记录 297619886}、
\href{https://www.luogu.com.cn/record/297620263}{传统版记录 297620263}。
单点最大用时分别为 64 ms、62 ms，内存峰值为 2.85 MB、2.40 MB。
在线测试覆盖无向与根为 1 的外向树；其他根、内向树和负权仍使用本地证据。
定理背景可参考 \href{https://oi-wiki.org/graph/matrix-tree/}{OI Wiki 矩阵树定理}。


\clearpage
\section{任意正模数的矩阵树计数}\label{knowledge-matrix-tree-mod}
实现见第~\ref{compact-MatrixTreeMod}~节（第~\pageref{compact-MatrixTreeMod}~页）。
拉普拉斯余子式表示生成树权值和的恒等式本身不要求模数为素数；
需要换掉的是依赖模逆元的行列式算法。
本节沿用第~\ref{knowledge-matrix-tree-weighted}~节（第~\pageref{knowledge-matrix-tree-weighted}~页）的三种构造，但使用欧几里得行列式消元。
输入边权和正模数均支持 signed 64 位，接口名为
\texttt{MatrixTreeMod::count}。行列式实现见第~\ref{compact-determinant_mod}~节（第~\pageref{compact-determinant_mod}~页），时间 $O(n^3\log(m+1)+|E|)$、空间 $O(n^2)$；此处 $m$ 为模数。

先把边权正规化到 $[0,m)$，每次矩阵元素加减仍先扩展至 128 位，
再取模存回 64 位。重边的权值和可能超过 64 位范围，不能先累加到普通整数中。
只构造删除根后的余子式，并将其存储直接交给行列式函数。
单点返回 $1\bmod m$，模 1 返回零。
负权与零权仍按代数权值处理，返回零不能直接解释为没有生成树。

独立参考程序逐个枚举边集，用任意精度整数计算权值和。
已检查多个合数模数、所有根与三种方向、断图、重边自环及 64 位极值权重；
另以大权值完全图的 Cayley 公式验证 160 点规模。
普通与 ASan/UBSan 均通过，尚无该接口的在线 AC。
此前 P6178 的 AC 属于素数模版本，不自动转移到本接口。
